\section{Conclusion}

We benchmarked classical, ant-colony and Physarum routing on the real
directed road network of central Bengaluru under microsimulated congestion,
measuring recovery from seven classes of road closure as well as path
quality, over 5{,}520 runs in a fully paired design.

The headline claim made for bio-inspired routing --- that retained internal
state buys recovery without recomputation --- holds, and holds in its
strongest form: warm start lets Physarum recover in a single iteration on
48.3\% of closures, which cold start never achieves, at an unchanged
per-iteration cost. But it is a trade and not a gift. Warm start is
significantly worse on the reported post-disruption path, and it degrades the
model's own extraction further, so reporting only the speed-up describes half
the effect.

The finding we think matters most is methodological. A Physarum model does
not emit a route; an extraction step does, and on a sparse directed network
that step decides the answer. Pooled over extraction methods our Physarum
matches the optimum on 90\% of undisrupted pairs; using only its own
maximum-flux walk, 23.3\%; and on warm recoveries it returns plain
Dijkstra's path two-thirds of the time. Because the median stretch is
exactly 1.000 inside every provenance class, a pooled figure cannot be
falsified. We argue that provenance should be reported as a matter of course
for this model class.

We also showed that adaptation improves extraction \emph{quality} while
reducing its \emph{reliability}, so the two must be measured together; that
no earlier stopping rule improves on the standard convergence criterion,
because the quality ceiling is already reached there; and that the sign of
the congestion effect on ACO belongs to the congestion model rather than to
the algorithm, which makes unqualified claims about behaviour ``under
congestion'' underdetermined.

Throughout, Dijkstra remained the honest baseline: an exact route in one
iteration and roughly 31\,ms, which no method here approached on wall-clock.
On a network of this size recomputation is the correct engineering choice,
and the case for adaptive methods must be made on regimes we did not test.

Future work should extend the benchmark across cities with different
topological character, pair the controlled closure protocol used here with
measured rather than synthetic demand, and evaluate the directed Physarum
formulation, whose behaviour on this substrate our results give specific
reason to want measured.

\section*{Reproducibility}

The complete experimental record --- the 5{,}520-row results file from which
every number in this paper derives, the network and congestion pipeline, all
seven figure-generating scripts, the hyperparameter sweeps, and the
pre-registered measurement protocols for the mechanism experiments --- is
available in the project repository. Every run carries its seed and a
configuration hash, and all figures regenerate from the committed results
file with a single command. Recorded failures are retained rather than
dropped, and a deliberately failing test marks the ACO exit criterion that
was missed.
